\documentclass[10pt,a4paper]{amsart}
\usepackage[margin=19mm]{geometry}
\usepackage{amsmath,amssymb,mathtools}
\usepackage[scr=rsfs,cal=euler]{mathalfa}
\usepackage{xfrac}
\usepackage[unicode,hidelinks,bookmarks=false]{hyperref}
\newcommand{\ie}{i.e.\ }
\newcommand{\cf}{cf.\ }
\newcommand{\resp}{respectively\ }
\newcommand{\Subsection}{\subsection}
\providecommand{\nobreakdash}{\nobreak\hbox{-}}
\setlength{\emergencystretch}{2em}
\multlinegap0pt
\usepackage{amssymb}
\usepackage{enumerate}
\usepackage[shortlabels]{enumitem}
\usepackage{caption}
\usepackage{xcolor}
\newtheorem{theorem}{Theorem}
\newtheorem{corollary}{Corollary}
\def\Z{{\mathbb Z}}
\title{Signed Clasp Numbers and Four-Genus Bounds}
\author{Charles Livingston}
\date{Source: arXiv:2206.10432v2 (CC BY 4.0)}
\begin{document}
\maketitle

\noindent\emph{Source and licence.} Signed Clasp Numbers and Four-Genus Bounds. Version arXiv:2206.10432v2 (CC BY 4.0), distributed by arXiv under the Creative Commons Attribution 4.0 International licence, http://creativecommons.org/licenses/by/4.0/, as recorded in the arXiv licence metadata for this exact version and shown on the abstract page https://arxiv.org/abs/2206.10432v2. Full licence notice: \texttt{licenses/arxiv-2206.10432-CC-BY-4.0.txt}.\par\smallskip

\noindent\textit{Editorial scope.} Counted unit: the article's theorem thm-boundscg (source0.tex lines 167--173). The article's complete original proof of it is delivered with it (source0.tex lines 174--180, one \texttt{proof} environment of 7 lines). No mathematics is written, repaired or re-derived: the statement and the proof are the article's own lines, in the article's own notation, and no retained line is substituted or re-flowed.  Retained uncounted as the source-local context the proof needs: lines 144--147; lines 163--164.  Nothing was omitted from any retained span, so the package carries no omission record.  No retained line contains a citation, so the wrapper carries no bibliography and no bibliography entry.  No retained line contains a figure environment, an image-inclusion command or drawing code, so no image is delivered and the package declares no asset dependency.  Editorial formatting only, in addition: an AMS \texttt{amsart} preamble that restates the article's own declarations and macros used by the retained lines, namely the environments \texttt{theorem}, \texttt{corollary} on separate counters, together with the standard packages those lines need.  The retained environments print in this document as Theorem 1, Corollary 1 in order of appearance, whereas the article prints the same items with the numbering of its own full document.  The complete native source of the article as distributed in the arXiv e-print for this exact version is bundled unchanged as \texttt{sources/127225/source0.tex}.
\begin{theorem}\label{thm-nonzero}
Let  $q = p^2$, where  $p$ is an odd prime.  Let $\chi$ denote a  character  that takes value  $e^{2 \pi i /p}$ on some generator of $H_1(M_2(B_q)) \cong \Z_q$.   Then
\[ \{\sigma(B_q, \chi^r)\}_{0<r< p}\ =\  \{  4r^2 - 2pr +1 \}_{0<r< p}.\]
\end{theorem}

\begin{corollary} \label{cor-zero} Suppose that $K$ bounds a genus one Seifert surface $F$,  $H_1(M_2(K)) \cong \Z_q$, and $V(\alpha, \alpha) = 0 $ for a simple closed curve  $\alpha $ representing a nontrivial homology class, $[\alpha] \in H_1(F)$.  Then $\sigma_1 \tau (K, \chi^0 ) =   0$ for all $\chi$.
\end{corollary}

\begin{theorem}\label{thm-boundscg} Assume $q = p^2$ where $p \ge 5$ is an odd   prime. % Let $\chi$ be a generator of the group of order $p$ characters on $H_1(M_2(B_q))$.
\begin{itemize}

\item There exists a  generator  $\chi$ of the group of order $p$ characters on $H_1(M_2(B_q))$   such that $\sigma_1 \tau(B_q, \chi) \le   \frac{9 - p^2}{4}$.
\item  $\sigma_1 \tau (B_q,  \chi^r) \le 0 $ for all $r$.
\end{itemize}
\end{theorem}

\begin{proof}  We consider the function   $f(r)  = 4r^2 - 2pr +1$ that appears in Theorem~\ref{thm-nonzero}  as a real quadratic in the variable $r$.
Its minimum occurs at $p/4$.
The closest integer point to $p/4$ is either $(p-1)/4$ or $(p+1)/4 $ depending on whether $p \equiv 1 \mod 4$ or $  p \equiv  3 \mod4$.   In both cases the value at this  point is $(5 - p^2)/4 < -1$.  Since $\sigma(B_q, \chi^r)$ and $\sigma_1\tau(B_q, \chi^r)$ differ by at most one, we have the first statement.

For integers $r$ with  $ 1 \le r \le p/2$, the maximum value of the quadratic $f(r)$ must be at an endpoint, either $r=1$ or $r = (p-1)/2$.  We compute $f(1) = (5 - 2p)$ and $f((p-1)/2) = 2 - p$. The larger of the two is $2 -p <1$.  Even upon adding 1, this is negative.  Thus, if $\sigma_1 \tau(B_q, \chi^r)$ were to be positive for some $r$, it would have to be at $r = 0$, where the value was shown to be 0 in Corollary~\ref{cor-zero}.

\end{proof}

\end{document}
